\documentclass[12pt]{article}

\usepackage[a4paper, total={6in, 8in}]{geometry}
\usepackage{textcomp}

\begin{document}
\setlength{\parindent}{0pt}

\def \t {\textrightarrow}
\def \v {\vspace{1em}}
\def \bi {\begin{itemize}}
\def \ei {\end{itemize}}
\def \s[#1] {\section*{#1}}
\def \ss[#1] {\subsection*{#1}}
\def \sss[#1] {\subsubsection*{#1}}

\s[Tacito - Germania]
Lui ha avuto esperienza diretta sul campo in germania \t il giudizio di tacito verso i barbari è un giudizio di valorizzazione e di rispetto \t è una popolazione pura \\
Nella loro forma mentis vivono relazioni prive di corruzione \t è un modello di cultura primitiva ma onesta \\
L'esibizione del loro nudo non è da vedersi come espressione maliziosa, ma espressione della loro socialità \t si lavavano tutti insieme \\
Sono uomini primitivi \t Adamo ed Eva sono nudi, ma quando escono dall'Eden si vergognano della loro nudità \\
Ma qual è il termine di confronto? \t anche Cesare nel suo De Bello Gallico parla di queste popolazioni \t rivisitazione di quelle popolazioni, ma senza lo scrupolo tacitiano di comporre un'opera etnogeografica/geoetnografica (SCRIVERE ENTRAMBI) \\
Nella Germania di Tacito l'osservazione è diretta \t Cesare invece narra, e cerca impersonalità, ma da lo stesso connotazione giudicante dei comportamenti di queste popolazioni \\
Ma l'ambiente era diverso \t in quel periodo, queste popolazioni minavano la centralità di Roma \t mentre Tacito sosta in Germania, e osserva direttamente \\
Tacito capisce che tutto cio che popolera le pagine degli annali (ovvero la corruzione, la violenza, lo sperpero di denaro, la impudicizia) non trovano invece riscontro nel mondo dei germani \t la Germania viene vista con un ribaltamento \\
Il carattere identitario barbarico viene ad essere quello primario di civiltà \t rispetto a Roma, l'ambiente era primitivo ma puro \\
Assistiamo a un ribaltamento (come manzoni con storia e antistoria)
\end{document}
